%!TeX root=00_Abschlussarbeit.tex

\selectlanguage{english}
\renewcommand{\abstractname}{Abstract}
%\begin{abstract}


\chapter*{Abstract - Noch nicht überarbeitet, Erste Fassung}
\addcontentsline{toc}{chapter}{Abstract}
\noindent
% TEXT ENGLISCH %
This thesis describes the development of a visual localization system for unmanned aerial vehicles based on fiducial markers in indoor environments.
The Robot Operating System 2 is used as the middleware framework; AprilTags are used as the marker framework.
The system uses the RGB images from two cameras as its input data source—one facing forward and one facing downward.
The image processing and navigation calculations are performed by a mini-computer mounted on the aircraft.
Position determination is based on the mapping of fiducial markers using simultaneous localization and mapping.
For the extrinsic calibration of the cameras, a custom-developed algorithm for determining the respective transformation matrix is described.
Position determination itself is achieved by detecting the previously mapped markers and calculating the poses of the camera lenses.
The position data from both sensors are processed, fused into a single pose, and transmitted to the flight controller.
The overall system enables GPS-independent positioning designed for indoor use.
Finally, a prototype application is presented that can be used to control the aircraft.

%\end{abstract}


\selectlanguage{ngerman}
\renewcommand{\abstractname}{Kurzfassung}
%\begin{abstract}
\chapter*{Kurzfassung - Noch nicht überarbeitet, Erste Fassung}
\addcontentsline{toc}{chapter}{Kurzfassung}
\noindent
% TEXT DEUTSCH %
Diese Arbeit beschreibt die Entwicklung eines visuellen Lokalisierungssystems für unbemannte Luftfahrzeuge auf Basis von Fiducial-Markern im Innenraum.
Als Middleware-Framework wird das Robot Operating System 2 verwendet; als Marker-Framework kommen AprilTags zum Einsatz.
Das System nutzt die RGB-Bilder zweier Kameras als Eingangsdatenquelle – eine frontal und eine nach unten ausgerichtet.
Die Berechnungen zur Bildverarbeitung sowie zur Navigation werden durch einen am Luftfahrzeug montierten mini-Computer durchgeführt.
Grundlage der Positionsbestimmung bildet die Kartierung der Fiducial-Marker mittels Simultaner Lokalisierung und Kartierung.
Zur extrinsischen Kalibrierung der Kameras wird ein eigens entwickelter Algorithmus zur Bestimmung der jeweiligen Transformationsmatrix beschrieben.
Die Positionsbestimmung selbst erfolgt durch die Erkennung der zuvor kartierten Marker sowie die Berechnung der Posen der Kameraoptiken.
Die Positionsdaten beider Sensoren werden aufbereitet, zu einer einzigen Pose fusioniert und an den Flugcontroller übermittelt.
Das Gesamtsystem ermöglicht eine GPS-unabhängige Lokalisierung, die für den Einsatz in Innenräumen ausgelegt ist.
Abschließend wird eine prototypische Anwendung vorgestellt, die zur Steuerung des Luftfahrzeugs eingesetzt werden kann.


%\end{abstract}